\section{Discussion}
\label{sec:discussion}

\textbf{What simulated drives over measured profiles provide.}
Ranking segments within a route turned out to be the easy part, which an untrained, speed-normalised vibration RMS does nearly as well as the best trained estimators.
Simulation supplied the mapping from vibration to IRI up to a level that depended on two simulator choices, the sensor path, whose effect reached the estimator through the features carrying vibration above 10\,Hz, and the vehicle population; no variant was built to reproduce the absolute gain of this car's vibration path, and the reference simulator's matched level was a fortunate combination of these choices.

\textbf{A blueprint for road agencies.}
For a new car type the results suggest three steps.
(i) Simulate the car over corridors whose profiles are already surveyed, with randomised vehicles and a sensor model; the ranking of segments will transfer.
(ii) Collect about 5\,km of paired drives over profiled road; one pass per segment sufficed here.
(iii) Train jointly on simulated and paired drives, which beat real-only training on average for every simulator variant and every $k$ we tested, though not on every route and draw; a level offset is slightly more accurate when the simulated level error is common to both road classes, which the paired drives reveal by class and a comparison of real and simulated spectra on the calibration roads anticipates (an association found after the fact).

\textbf{How to evaluate.}
Because a within-route Spearman correlation is unaffected by calibration, it cannot separate a trained estimator from an untrained statistic; IRI estimators should be judged on held-out routes by absolute error and level per road class.
Segment errors are spatially correlated: 5\,km-block intervals of paired differences were \infWidthRatioMin--\infWidthRatioMax{} times as wide as segment-level ones, and route-level intervals with four routes include zero for most comparisons (Table~\ref{S-tab:s_inference}), so the supported differences describe these four routes, not a population of routes, and benchmarks on road data should resample spatial blocks or routes, not segments.
Smooth roads are driven fast, so estimators trained on real cars can use speed as a cue for low roughness, consistent with their over-correction on a fixed road; we recommend reporting such fixed-road elasticities from repeated passes, with the within-pass speed variability as a control, instead of speed sweeps in simulation alone.
All errors are relative to one P79 survey (a second device: Table~\ref{S-tab:s_aran}).

\section{Limitations}
\label{sec:limitations}
\emph{External validity.}
LiRA-CD contains one car model with one type of telematics unit on four routes in one city region; road class and route are confounded, and with four routes no statement about route-to-route variation is inferential.
Whether the simulated mapping transfers up to a level, and how much paired data the level needs, for other cars, sensors, speeds and road networks is unknown.
Transfer to other vehicles and to phones in holders was examined only in simulation (Supplement~\ref{S-sec:s_controlled}).
\emph{Exploratory analyses.}
Except for the comparison of the reference simulator with real training, all real-data results come from analyses designed after the test routes had been evaluated; they generate hypotheses rather than confirm them.
\emph{Combination methods.}
We compared three ways of combining simulated and real data with one learner and fixed settings; network fine-tuning or domain adaptation might rank differently.
\emph{Simulator.}
The simulator is a linear half-car without roll or body flexibility, driven by the mean of the two wheel-path profiles, with parameter ranges for sensitivity rather than a calibrated fleet; its vibration differs from the real car's at 2--5\,Hz and above 10\,Hz and in its speed scaling.
\emph{Other.}
M7 is not a reproduction of~\cite{skar2024inversion} and the CNN was not tuned; the reference surveys preceded the car passes by two to fifteen months, segments are 100\,m long, and we did not train on real data from several cities.
